\chapter{PERANCANGAN PROTOTIPE/PRODUK}

\section{Deskripsi Umum}

% TODO: Jelaskan gambaran besar produk yang dirancang.

Sistem yang dirancang dalam penelitian ini adalah lorem ipsum dolor sit amet, consectetur adipiscing elit. Sed do eiusmod tempor incididunt ut labore et dolore magna aliqua.

% TODO: Ganti dengan diagram arsitektur sistem Anda.
% \begin{figure}[H]
%   \centering
%   \includegraphics[width=0.9\textwidth]{figures/arsitektur-sistem}
%   \caption{Diagram Arsitektur Sistem}
%   \label{fig:arsitektur-skripsi}
% \end{figure}

\section{Alur Pengembangan Produk}

% TODO: Jelaskan metodologi pengembangan produk.

Pengembangan sistem dilakukan dengan menggunakan pendekatan Research and Development (R\&D). Metodologi ini dipilih karena penelitian berfokus pada pengembangan produk inovatif.

\section{Perancangan Arsitektur Sistem}

\subsection{Komponen Utama}

% TODO: Jelaskan komponen-komponen utama sistem.

Arsitektur sistem terdiri dari komponen-komponen berikut:

\begin{itemize}
    \item \textbf{Komponen A:} Lorem ipsum dolor sit amet, consectetur adipiscing elit.
    \item \textbf{Komponen B:} Sed do eiusmod tempor incididunt ut labore et dolore magna aliqua.
    \item \textbf{Komponen C:} Ut enim ad minim veniam, quis nostrud exercitation ullamco laboris.
\end{itemize}

\subsection{Alur Kerja Sistem}

% TODO: Jelaskan alur kerja sistem secara detail.

Lorem ipsum dolor sit amet, consectetur adipiscing elit. Sed do eiusmod tempor incididunt ut labore et dolore magna aliqua. Ut enim ad minim veniam, quis nostrud exercitation ullamco laboris nisi ut aliquip ex ea commodo consequat.

\section{Perancangan Keamanan / Aspek Teknis Lain}

% TODO: Sesuaikan dengan aspek teknis produk Anda.

\subsection{Aspek Keamanan}

Lorem ipsum dolor sit amet, consectetur adipiscing elit. Sed do eiusmod tempor incididunt ut labore et dolore magna aliqua.

\subsection{Aspek Lainnya}

Lorem ipsum dolor sit amet, consectetur adipiscing elit. Sed do eiusmod tempor incididunt ut labore et dolore magna aliqua.
